\section{Transporting the pole coefficient through the block bridge}
\label{sec:ex-quadratic-pole-transport}

The annulus calculation concerns a fine-lattice determinant, whereas
the stationary output uses its pullback to coarse fields. We prove
that their leading quadratic coefficients agree. This uses both
quadratic matching equations for the selected bridge, including its
normalization at zero momentum.

Let $\Pi_f$ be a common quadratic picture for the fine pole functional
\[
 P_f(x)=\frac12\log\det\left(K(x)+d(x)d(x)^*
                         +\sum_a s_a^*C(x)s_a\right)
                      -\log\det\left(d(x)^*d(x)+C(x)\right).
\]
The notation is that of
Section~\ref{sec:ex-stationary-determinant-prices}. Suppose its
finite-volume recording error is at most the fine site cardinality
times $e_m$, with $e_m\to0$ along the selected periods. Let $\Pi_c$
be the common picture of the same functions pulled back by the
installed linear bridge $W$, with coarse per-site errors tending
to zero. Both pictures are holomorphic on positive strips. Fix the
block side $L$ throughout the argument.

The standing data are those of one stationary step, and all of them
refer to the same step. The gauge group is simple. The coarse sides
$m_i$ tend to infinity and the fine sides are $n_i=Lm_i$. The path $R$
has path bounds in the clip $p$ and is attached to the fine bed $B$;
the metric row of $B$ has a strip bound at width $\sigma/L$, the
floor of $B,p$ holds with one constant $\varepsilon>0$ at every torus
side, and the sworn condition of $B$ holds at every $n_i$. The metric
row of the coarse bed $T$ has a strip bound at width $\sigma$. The
bridge $\Phi$ is holomorphic and periodic on the strip of width
$\sigma$ and is the common Fourier representation of the installed
matrices $W_i$, which satisfy both matching equations
\eqref{eq:inherited-matching} for $B$ and $T$, with block mean $Q_i$
and penalty coefficient $\kappa$. The frame and kinetic response data
of Section~\ref{sec:ex-selected-kinetic-response}, from which the rank
response is built, are formed with these same $W_i$. The pole
functional $P_f$ is formed with the regulator $\mu=L^{-1}$, and the
fine pole picture $\Pi_f$ records it at the sides $n_i$. The argument
of this section does not use the stationary branch or its Hessian gap,
although in the source they belong to the same packaged data of the
step.

The block determinant matrix $M_{\rm rank}$ is the direct sum of the
fine pole matrix and an identity on coarse scalar coordinates,
after the coordinate permutation separating the two kinds of
variables. Thus its determinant equals the fine pole determinant.
The scalar determinant is the same in both presentations. The
Fourier matrix identities proved above therefore identify the
coarse rank picture with the picture of $P_f\circ W$; this is
an equality for the actual selected finite-volume functions.

The pole function is invariant under local gauge changes, constant
color rotations and unsigned coordinate permutations. For gauge changes this
follows by conjugating every factor: the fine Hessian in left
coordinates, the covariant difference, its resolvent regulator and
the component maps all transform by the corresponding orthogonal
adjoint actions. Determinants are unchanged. The same argument,
with coordinate permutations whose direction signs are all positive,
proves this permutation invariance. Its first derivative at zero vanishes by the
simple-color argument in
Section~\ref{sec:ex-thermal-contact-cancellation}.
Passing the finite-volume identities to the common picture, as in
Section~\ref{sec:stationary-output-construction}, shows that its
bilinear matrix $S(p)$ has both Ward identities
\[
 S(p)d_0(p)=0,\qquad d_0(-p)^T S(p)=0,
\]
The first-order part of the Ward identities forces $S(0)=0$.
To prove that the linear part also vanishes, fix the two color indices
and write it as $S_{ab}(p)=\sum_c T_{ab,c}p_c+O(|p|^2)$.
The two Ward identities, at quadratic order, give
\[
 T_{ab,b}=T_{ab,a}=0,\qquad T_{ab,c}+T_{ac,b}=0.
\]
These equations also kill the case $a=b$: then
$T_{aa,c}=-T_{ac,a}=0$. If $a,b,c$ are distinct, the unsigned
permutation interchanging $b,c$ and fixing $a$ gives
$T_{ab,c}=T_{ac,b}$; the displayed antisymmetry makes both zero.
All repeated-index cases are already zero. Thus $S$ vanishes to
order at least two. This uses only unsigned permutations; no
reflection symmetry of the pole functional is assumed. The
finite-volume errors disappear after division by the site
cardinality. No rate uniform in $L$ is used in this step.

Write $H(p)$ for the bilinear matrix of $\Pi_c$ and $N_l(p)$
for the fine alias amplitudes of $W$. The finite-wave expansion
from Section~\ref{sec:inherited-loop-transport} gives
\begin{equation}\label{eq:ex-pole-alias-pullback}
 H(p)=L^4\sum_l N_{-l}(-p)^T S(k_l(p))N_l(p),
                  \qquad k_l(p)=\mu(p+2\pi l).
\end{equation}
We spell out the role of errors here. A pulled-back coordinate wave
is a finite sum of fine alias waves with amplitudes $N_l(p)$.
For fixed $L$, these amplitudes remain bounded as the real momentum
approaches a specified value. Multilinearity bounds the error on
each pair by the original per-site error times the finite sum of
the products of its amplitudes. The ratio of fine to coarse site
cardinalities is $L^4$. Hence the error still tends to zero.
Opposite aliases are the only surviving pairs, by fine translation
invariance; all other pairs tend to zero. Continuity and a dense
sequence of external grid momenta prove
\eqref{eq:ex-pole-alias-pullback} at every real momentum.

For $l\ne0$, the transverse projection of $N_l(p)$ has order two
at $p=0$, by the all-period kinetic floor and the primal matching
equation proved in
Section~\ref{sec:inherited-loop-transport}. The two Ward
identities of $S$ remove the longitudinal parts of both amplitudes.
Thus every nonprincipal term of
\eqref{eq:ex-pole-alias-pullback} has order at least four. This
argument does not require $S$ to vanish at the nonprincipal alias.

It remains to determine the principal amplitude at zero. The dual
matching equation, the second equation of
\eqref{eq:inherited-matching}, passes from the grids to all real
momenta by the density argument already used. With $Q(p)$ the symbol
of the block mean it reads
\[
                  Q(p)\Phi(p)=I-\kappa^{-1}K_c(p).
\]
Since $K_c(0)=0$, its value at zero is $Q(0)\Phi(0)=I$.
Each nonprincipal alias amplitude at zero is a pure gauge vector,
because its transverse projection vanishes there. The block mean
commutes with the gauge difference, so it maps this vector to a
coarse gradient at momentum zero, which is zero. Hence those
aliases contribute nothing to $Q(0)\Phi(0)$. On the principal
constant alias, the normalized transverse weights sum to one and
the straight block path has $L$ edges; therefore the mean acts
as multiplication by $L$. We conclude
\begin{equation}\label{eq:ex-principal-amplitude}
                              N_0(0)=\mu I.
\end{equation}
In particular this normalization is a consequence of the actual
two matching equations, rather than a choice of Fourier convention
made after the comparison.

Analyticity gives $N_0(p)=\mu I+O(|p|)$. Since $S$ vanishes
to order two and $k_0(p)=\mu p$, expanding the two principal
amplitudes yields
\[
             H(p)-L^4\mu^2S(\mu p)=O(|p|^3)
\]
along every ray. The nonprincipal terms have already been shown
to be of order four. If $S_2$ is the homogeneous quadratic term
of $S$, the quadratic term of the expression on the right is
\[
                        L^4\mu^2\mu^2S_2(p)=S_2(p).
\]
Thus $H$ and $S$ have identical quadratic Taylor terms, and in
particular identical drift coefficients in every color component.

We state the result and its use.

\begin{proposition}[Transport of the pole coefficient]
\label{prop:ex-pole-coefficient-transport}
Under the standing data listed at the beginning of this section, let
$\Pi_f$ be a two-leg picture, with components holomorphic, periodic and
bounded on a strip of positive width, which records the fine pole
functional $P_f$ of $p,R$ with regulator $L^{-1}$ at every fine side
$n_i$ with error $n_i^4e_i$, where $e_i\to0$. Let $\Pi_c$ be the
picture associated, as in Lemma~\ref{lem:ex-pole-recording}, with the
rank response kernel, and assume that this kernel is jointly
holomorphic, periodic and bounded on a strip of positive width. Then
for every color $c$
\[
 \alpha_c(\Pi_c)=\alpha_c(\Pi_f).
\]
\end{proposition}
\begin{proof}
By the block identity for the determinant matrix recalled above and
Lemma~\ref{lem:ex-pole-recording}, $\Pi_c$ records $P_f\circ W_i$ on
the $i$th coarse torus with error equal to the coarse site count times
a quantity tending to zero. The argument of this section then shows
that the bilinear matrices $H$ of $\Pi_c$ and $S$ of $\Pi_f$ have the
same quadratic homogeneous term at zero momentum. The coefficient
$\alpha_c$ is twice the real part of one entry of that term.
\end{proof}

In the one-step construction three coefficients occur, and the
proposition relates two of them. Let $\mathcal Y_{\rm det}$ be the
picture of the completed determinant correction
$\mathcal B_t$ of Section~\ref{sec:stationary-germ-symmetry} along the
stationary branch, with regulator $\mu=L^{-1}$. The scalar of the
output path is $\alpha=\alpha_c(\mathcal Y_{\rm det})$
(Proposition~\ref{prop:stationary-output-path}). The inherited loop
does not enter $\alpha$: by the moving Haar identity
\eqref{eq:ex-thermal-determinant} the thermal function is the
inherited loop plus $\mathcal B_t$ along the same branch, and the
inherited picture has order four
(Proposition~\ref{prop:inherited-loop-picture}), so its quadratic
coefficient is zero. The comparison of the completed determinant
response with the rank response is the subject of
Sections~\ref{sec:ex-stationary-determinant-prices}
and~\ref{sec:ex-thermal-contact-cancellation}: it is to supply a
constant $C_{\rm cmp}$, fixed before $L$, with
\[
 |\alpha_c(\mathcal Y_{\rm det})-\alpha_c(\Pi_c)|\leq C_{\rm cmp}.
\]
This is an input here. Finally
Proposition~\ref{prop:ex-annulus-estimate} supplies a fine pole picture
$\Pi_f$ for the same clip, path, fine bed and regulator, recorded at
the same fine sides $n_i$, with
\[
       |\alpha_c(\Pi_f)+\lambda_G\log L|\leq D_0 .
\]
The three statements concern the same data, and the triangle
inequality gives
\[
       |\alpha+\lambda_G\log L|
       \leq D_0+C_{\rm cmp}.
\]
The hypotheses \eqref{eq:ex-annulus-hypotheses} of the annulus
estimate, for the fine bed and the fine sides, are among the standing
data of this section together with the bound on the direct coefficient
of $p$ and the strip bound of $B$ at a fixed width.
